\section{Coefficient recurrences and support alternatives}
\label{h:ev:section}

The main block estimate propagated a scale inside a fixed radius
power. An equivalent formulation propagates the coefficient of that
power. Keeping the exponent variable makes the different terminal
accuracy powers visible, while one recurrence handles all of them.
The density and exact ellipsoid estimates have already been proved.
We first record their spherical specialization and two optional
calculations, then give the common coefficient recurrence.

Write \(\sigma=\sigma_d\) and \(n=d-1\). For unit vectors \(u,v\),
put \(\operatorname{angle}(u,v)=\arccos\langle u,v\rangle\).
The following specialization keeps the local exponent variable while
fixing a common moment order.

\begin{corollary}[The dimension-scale parameter regime]
Under the hypotheses of Proposition~\ref{h:dv:ambient-gram}, suppose that
\begin{equation}
 d\ge16,\qquad p=\lfloor d/8\rfloor,\qquad
 1\le m\le p,\qquad \frac{d-1}{2}\le\beta\le\frac d2,
 \label{h:dv:parameters}
\end{equation}
Then, for every \(0\le j\le p-2\),
\begin{equation}
 \beta-j-m\ge\frac{d-1}{2}-2p+2
 \ge\frac d4+\frac32,
 \qquad J_j\le C^d,
 \label{h:dv:dimension-margin}
\end{equation}
where \(C\ge1\) is an absolute constant.  Thus the right side of
\eqref{h:dv:gram-product} is at most
\(C^{d(p-1)}H^pR^{\beta p-m(p-1)}\), uniformly also for short blocks
\(m<p\).
\end{corollary}

\begin{proof}
Under \eqref{h:dv:parameters}, \(p\le d/8\) gives the displayed
margin.  Also \(m/(\beta-j-m)\le1/2\), while
\(3^j2^\beta\le3^{d/8}2^{d/2}\).  Enlarging one absolute \(C\) gives
\(J_j\le C^d\), and multiplying over the \(p-1\) values of \(j\)
proves the specialization.
\end{proof}


Here \(p\ge d/16\) is the moment order and maximum block length,
whereas \(m\) is the current number of rows. The constants below
are absolute unless a dependence is stated.

\subsection{Spherical restrictions and a second moment calculation}

\begin{lemma}[Exact projection density from an ambient ball bound]
\label{h:ev:density}
Assume \eqref{h:dv:parameters}, and let \(H,R>0\) and
\(z\in\mathbb R^d\).  Let \(f:S^{d-1}\to[0,1]\) be Borel
measurable, let \(\nu=f\sigma\), and suppose
\(\nu(B(x,r))\le H r^\beta\) for every \(x\in\mathbb R^d\) and \(r>0\).
For \(\mu=\mathbf 1_{B(z,R)}\nu\), the pushforward \(A_\#\mu\) has,
for almost every \(A\), a Lebesgue density \(F(A,\cdot)\).  There is
an absolute \(C_{\mathrm D}\ge1\) for which a jointly measurable
nonnegative choice of these densities satisfies
\begin{equation}\label{h:ev:density-moment}
 \int F(A,y)^p\,d\lambda_m(A,y)
 \le C_{\mathrm D}^{d(p-1)}
       H^pR^{\beta p-m(p-1)}.
\end{equation}
The density represents the exact labels \(As\).
\end{lemma}

\begin{proof}
The restriction \(\mu\) is supported on \(B(z,R)\) and inherits the
all-radius bound. Theorem~\ref{h:dv:finite-density}, together with
\eqref{h:dv:dimension-margin}, gives its strong joint \(L^p\) density
and \eqref{h:ev:density-moment}. The fiber identification in
Lemma~\ref{h:dv:joint-density} gives the asserted common almost-everywhere
choice of projection densities.
\end{proof}

The same bound can be obtained from ball averages with a less sharp
explicit prefactor. This calculation is independent of the choice of
terminal coefficient or support domain.

For \(\delta>0\), set
\[
 F^{\mathrm{ball}}_\delta(A,y)
 =\frac{1}{v_m\delta^m}
   \int\mathbf 1_{\{\|As-y\|\le\delta\}}\,d\mu(s).
\]
These functions are jointly measurable.  For almost every \(A\) they
are the ball averages of the exact density \(F(A,\cdot)\); hence
Lebesgue differentiation and Fubini give
\(F^{\mathrm{ball}}_\delta\to F\) for \(\lambda_m\)-almost every
\((A,y)\) along \(\delta=1/k\).  In particular this limit, set to zero
where it does not exist finitely, is a jointly measurable density
version.  
Expand the \(p\)th moment and apply Tonelli's theorem.  The intersection
of the \(p\) balls of radius \(\delta\) about \(As_1,\ldots,As_p\) has
volume at most \(v_m\delta^m\), and is empty unless
\(\|A(s_i-s_1)\|\le2\delta\) for every \(i\ge2\).  Therefore
\begin{align}
 \int (F^{\mathrm{ball}}_\delta)^p\,d\lambda_m
 &\le (v_m\delta^m)^{-(p-1)}
 \int
 \Pr_A\{\|A(s_i-s_1)\|\le2\delta,\ 2\le i\le p\}
 \,d\mu^{\otimes p}.
 \label{h:ev:ball-intersection}
\end{align}
For a full-column-rank \(V=[s_2-s_1,\ldots,s_p-s_1]\), the rows of
\(AV\) are independent centered Gaussian vectors with covariance
\(V^{\mathsf T}V\).  Their joint density is bounded above by
\[
 (2\pi)^{-m(p-1)/2}\det(V^{\mathsf T}V)^{-m/2}.
\]
In the difference variables \(AV\), the event in
\eqref{h:ev:ball-intersection} lies in the product of \(p-1\)
radius-\(2\delta\) balls in \(\mathbb R^m\), whose
\(m(p-1)\)-dimensional volume is \((v_m(2\delta)^m)^{p-1}\).
The exceptional rank-deficient tuples are null by
\eqref{h:dv:gram-product}.  It follows that
\begin{equation}\label{h:ev:ball-moment}
 \int (F^{\mathrm{ball}}_\delta)^p\,d\lambda_m
 \le 2^{m(p-1)}(2\pi)^{-m(p-1)/2}
     \int\det(V^{\mathsf T}V)^{-m/2}\,d\mu^{\otimes p}.
\end{equation}
All factors \(v_m\delta^m\) have canceled.  Combining
\eqref{h:ev:ball-moment} with \eqref{h:dv:gram-product}, absorbing
\(2^{m(p-1)}\) into an absolute constant to the power \(d(p-1)\),
and using Fatou's lemma proves \eqref{h:ev:density-moment}.

For the full-block choice
\(\beta=(d-1)/2\) and \(m=p=\lfloor d/8\rfloor\),
\eqref{h:ev:ball-moment} retains the factor \(2^{p(p-1)}\).
The margin in \eqref{h:dv:dimension-margin} is in particular at least \(d/8\).
Taking the \(p\)th root gives the form
\begin{equation}\label{h:ev:half-sphere-density}
 \|F\|_{L^p(\lambda_p)}
 \le C_{\mathrm{ball}}^d H R^{\beta-p(1-1/p)}.
\end{equation}

For Gaussian averages no new identity is needed:
\eqref{h:proj:gaussian-identity} with \(k=p\) and
\eqref{h:dv:gaussian-moment-bound} give the exact covariance
\(V^{\mathsf T}V+\delta^2(I_{p-1}+\mathbf1\mathbf1^{\mathsf T})\)
and prefactor \((2\pi)^{-m(p-1)/2}\).
The compact-test, joint-\(L^1\), and pointwise passages all give the
same moment estimate under their respective hypotheses. In particular,
setting \(\beta=d/2\) in \eqref{h:ev:density-moment} gives, for every
\(1\le m\le p\),
\begin{equation}\label{h:ev:half-ambient-density}
 \mathbb E_A\int F(A,y)^p\,dy
 \le C_{\mathrm D}^{d(p-1)}
       H^pR^{(d/2)p-m(p-1)}.
\end{equation}

\subsection{Enclosing balls as an alternative support}

The exact image domain \(\mathcal E_{z,R}\) has the volume in
Lemma~\ref{h:ev:ellipsoid}. The larger operator-norm ball
\(\mathcal S_{z,R}\) from \eqref{h:dv:support-sets} also suffices;
the following proof records its uniformity for short blocks.

\begin{lemma}[Projected support volume]\label{h:proj:volume}
For a standard Gaussian \(m\)-by-\(d\) matrix \(A\), where
\(1\le m\le d\),
\begin{equation}\label{h:proj:volume-bound}
 \E\,\vol_m B(0,\|A\|_{\op})\le\exp(Cd).
\end{equation}
\end{lemma}

\begin{proof}
Choose \(1/4\)-nets of the unit spheres in \(\R^d\) and \(\R^m\)
with at most \(9^d\) and \(9^m\) points. Approximating both arguments
in the bilinear supremum loses at most \(\|A\|_{\op}/2\), so
\[
 \|A\|_{\op}\le2\max_{u,v}|v^{\mathsf T}Au|
\]
over the net pairs. Each scalar is standard Gaussian. For such a
scalar \(Z\),
\(\E|Z|^m\le(C\sqrt m)^m\): use
\(|z|^m e^{-z^2/4}\le(2m/e)^{m/2}\) and
\(\E e^{Z^2/4}=\sqrt2\). Bounding the maximum's \(m\)-th power
by the sum gives
\[
 \E\|A\|_{\op}^m\le2^m9^{m+d}(C\sqrt m)^m.
\]
The unit \(m\)-ball volume \(v_m\) obeys
\(v_m\le(2\pi e/m)^{m/2}\): integrate \(e^{-\|y\|^2/2}\)
over the ball of radius \(\sqrt m\) and compare with its full Gaussian
integral. Multiplication cancels \(m^{m/2}\), proving
\eqref{h:proj:volume-bound} uniformly even for a short block.
\end{proof}

Since
\(\lambda_m(\mathcal S_{z,R})
 =v_mR^m\mathbb E\|A\|_{\mathrm{op}}^m\), this proves
\begin{equation}\label{h:dv:enclosing-volume}
 \lambda_m(\mathcal S_{z,R})\le C_{\mathrm V}^dR^m
 \qquad(1\le m\le d).
\end{equation}

There is a second net calculation with a different explicit factor.
Maximal \(1/2\)-separated subsets of the unit spheres in
\(\mathbb R^d\) and \(\mathbb R^m\) are \(1/2\)-nets of sizes at most
\(5^d\) and \(5^m\).  Approximating the supremum of a linear form in
each argument in turn gives
\(\|A\|_{\mathrm{op}}\le4\max|v^{\mathsf T}Au|\) over the net pairs.
Each scalar in this maximum is a standard Gaussian \(Z\), and
\[
 \mathbb E|Z|^m
 \le\sqrt2(2m/e)^{m/2},
\]
because \(|t|^me^{-t^2/4}\le(2m/e)^{m/2}\) and
\(\mathbb E e^{Z^2/4}=\sqrt2\).  Bounding the maximum's \(m\)th power
by the sum therefore gives
\begin{equation}
 \mathbb E\|A\|_{\mathrm{op}}^m
 \le4^m5^{d+m}\sqrt2(2m/e)^{m/2}.
 \label{h:dv:half-net}
\end{equation}
Again multiplying by \(v_mR^m\) proves
\eqref{h:dv:enclosing-volume}, with a possibly larger absolute
constant.  This calculation retains the \(1/2\)-net factor \(4\).

\subsection{Terminal coefficients}

The geometric exponent \(\beta\) may be chosen before either of
the following terminal estimates. The hemisphere proof keeps the
explicit coefficient \(4^n\); the polar proof records two useful
denominator bounds.

\begin{lemma}[An ambient ball bound from a hemisphere]
\label{h:dv:hemisphere-cap}
Let \(d\ge2\), \(n=d-1\), and let \(\sigma_d\) be uniform probability
on \(S^{d-1}\).  Then
\begin{equation}
 \sigma_d(B(z,r))\le(4r)^n
 \qquad(z\in\mathbb R^d,\ r>0).
 \label{h:dv:hemisphere-ball}
\end{equation}
Let \(\epsilon>0\), let \(\kappa\) be any probability law on
\(S^{d-1}\), independent of the target, and put
\(h(s)=\int\mathbf 1_{\{\|s-u\|\le\epsilon\}}\,\kappa(du)\).
For \(0<\beta<n\), every \(z\in\mathbb R^d\), and every \(r>0\),
\begin{equation}
 \int_{B(z,r)}h\,d\sigma_d
 \le4^n\min\{r^n,\epsilon^n\}
 \le4^n\epsilon^{n-\beta}r^\beta.
 \label{h:dv:terminal-coefficient}
\end{equation}
The same upper bound holds for the angular-success function, since
angular distance at most \(\epsilon\) implies chordal distance at most
\(\epsilon\).
\end{lemma}

\begin{proof}
Suppose first that \(0<r\le1/4\) and that \(B(z,r)\) meets the sphere.
Choose a point \(a\) of the intersection and rotate it to a coordinate
axis.  Every other point of the intersection is within chordal distance
\(2r\) of \(a\), lies in the corresponding open hemisphere, and
projects into the radius-\(2r\) ball in \(\mathbb R^n\).  On that
projected ball the graph area element is at most \(2\), because
\(2r\le1/2\).  Thus the area of the patch is at most
\(2v_n(2r)^n\), where \(v_n\) is the volume of the unit \(n\)-ball.
The two hemisphere graphs together have area at least \(2v_n\).
The resulting probability is at most \((2r)^n\), and hence at most
\((4r)^n\). An empty intersection is harmless. For \(r>1/4\),
the bound \(\sigma_d(B(z,r))\le1<(4r)^n\) proves
\eqref{h:dv:hemisphere-ball}.

For fixed \(u\), the successful part of \(B(z,r)\) is contained both
in \(B(z,r)\) and in \(B(u,\epsilon)\).  Apply
\eqref{h:dv:hemisphere-ball} to each ball, take the smaller bound,
and integrate in \(\kappa\).  Finally
\(\min\{r^n,\epsilon^n\}\le\epsilon^{n-\beta}r^\beta\), by considering
\(r\le\epsilon\) and \(r\ge\epsilon\) separately.
\end{proof}

\begin{lemma}[Polar cap estimates]
\label{h:ev:polar-cap}
Let \(d\ge2\), \(n=d-1\), and let \(\sigma\) be uniform on \(S^{d-1}\).
There is an absolute \(C_0\ge1\) such that
\(\sigma(B(z,R))\le C_0^dR^n\) for every ambient center \(z\) and
every \(R>0\).  Let \(\epsilon>0\).  If a unit output is drawn from a
law independent of the signal, its success function \(h(s)\), for either angular error
at most \(\epsilon\) or chordal error at most \(\epsilon\), satisfies
\begin{equation}\label{h:ev:polar-terminal}
 \int_{B(z,R)}h\,d\sigma
 \le C_0^d\min(R^n,\epsilon^n)
 \le C_0^d\epsilon^{n-\beta}R^\beta
 \qquad(0<\beta<n).
\end{equation}
\end{lemma}

\begin{proof}
In spherical polar coordinates, the area element at polar angle \(t\)
is \(\sin^{d-2}t\,dt\) times the area element on \(S^{d-2}\).
Dividing the cap area by the total area therefore shows that a cap
of angular radius \(0\le\theta\le\pi\) has mass
\[
 \frac{\int_0^\theta\sin^{d-2}t\,dt}
      {\int_0^\pi\sin^{d-2}t\,dt}.
\]
The numerator is at most \(\theta^n/n\), since \(\sin t\le t\).
Restricting the denominator to two different intervals gives
\begin{align}
 \int_0^\pi\sin^{d-2}t\,dt
 &\ge \frac\pi3\left(\frac{\sqrt3}{2}\right)^{d-2}
 &&\text{from }[\pi/3,2\pi/3],
 \notag\\
 \int_0^\pi\sin^{d-2}t\,dt
 &\ge \frac{2\pi}{3}\,2^{-(d-2)}
 &&\text{from }[\pi/6,5\pi/6].
 \label{h:ev:polar-denominators}
\end{align}
Either estimate gives an upper bound \(C^d\theta^n\) for the cap.

If \(B(z,R)\) meets the sphere, choose \(s_0\) in the intersection.
Every other point in it has chordal distance at most \(2R\) from
\(s_0\).  The angle between two unit vectors is at most \(\pi/2\)
times their chordal distance.  Thus their angle is at most \(\pi R\)
when \(R\le1\); for \(R\le1/4\) the bound
\(2\arcsin R\le4R\) is also available.  The cap estimate and the
trivial mass bound one for \(R\ge1\) give the asserted ambient-ball
bound after enlarging \(C_0\).

For a fixed unit output, angular error at most \(\epsilon\) implies
chordal error at most \(\epsilon\), because
\(2\sin(\theta/2)\le\theta\).  The global success mass is therefore
at most \(C_0^d\epsilon^n\), while its mass in \(B(z,R)\) is at most
the ambient-ball bound \(C_0^dR^n\).  Averaging these bounds over
the output law proves the minimum in \eqref{h:ev:polar-terminal}.
If \(R\le\epsilon\), then \(R^n\le\epsilon^{n-\beta}R^\beta\);
if \(R\ge\epsilon\), then
\(\epsilon^n\le\epsilon^{n-\beta}R^\beta\).
\end{proof}

Lemma~\ref{h:dv:hemisphere-cap} gives the
alternative coefficient
\begin{equation}\label{h:ev:hemisphere-terminal}
 \int_{B(z,R)}h\,d\sigma
 \le4^n\min(R^n,\epsilon^n)
 \le4^n\epsilon^{n-\beta}R^\beta.
\end{equation}
For \(\beta=n/2\), the looser coefficient
\(L_0=4^d\epsilon^\beta\) follows immediately from
\eqref{h:ev:hemisphere-terminal}.  For \(\beta=d/2\),
\eqref{h:ev:polar-terminal} gives
\begin{equation}\label{h:ev:half-ambient-terminal}
 H_0=C_0^d\epsilon^{d-1-d/2}
     =C_0^d\epsilon^{d/2-1}.
\end{equation}

\subsection{One coefficient recurrence}

\begin{proposition}[Backward coefficient step]
\label{h:ev:coefficient-step}
Assume \eqref{h:dv:parameters}.  Let \(H>0\) and let \(N\ge1\) be an
integer.  Let
\(h_1,\ldots,h_N:S^{d-1}\to[0,1]\) be Borel functions satisfying
\[
 \int_{B(x,r)}h_w\,d\sigma\le H r^\beta
 \qquad(x\in\mathbb R^d,\ r>0,\ 1\le w\le N).
\]
Let \(g_w(A,y)\ge0\) be Borel functions with \(\sum_wg_w\le1\), and
define
\[
 h(s)=\mathbb E_A\sum_{w=1}^N g_w(A,As)h_w(s).
\]
There is an absolute \(C_{\mathrm{step}}\ge1\) such that
\begin{equation}\label{h:ev:coefficient-bound}
 \int_{B(z,R)}h\,d\sigma
 \le C_{\mathrm{step}}^dN^{1/p}H R^\beta
 \qquad(z\in\mathbb R^d,\ R>0).
\end{equation}
\end{proposition}

\begin{proof}
For the fixed ball \(B(z,R)\), take
\(\mu_w=\mathbf1_{B(z,R)}h_w\sigma\) in
Lemma~\ref{h:ev:density}.  Its exact density \(F_w\) is supported
on the common domain \(\mathcal E_{z,R}\) of
Lemma~\ref{h:ev:ellipsoid}.  The image identity gives
\[
 \int_{B(z,R)}h\,d\sigma
 =\int_{\mathcal E_{z,R}}\sum_w g_w(A,y)F_w(A,y)
       \,d\lambda_m(A,y).
\]
The density condition of Lemma~\ref{h:alt:selection} applies on this
domain with its norm parameter equal to
\[
 \mathcal U=C_{\mathrm D}^{d(1-1/p)}
             H R^{\beta-m(1-1/p)}.
\]
Its conclusion and Lemma~\ref{h:ev:ellipsoid} give
\begin{align*}
 \int_{B(z,R)}h\,d\sigma
 &\le N^{1/p}\mathcal U
       \lambda_m(\mathcal E_{z,R})^{1-1/p}\\
 &\le N^{1/p}C_{\mathrm D}^{d(1-1/p)}
       H R^{\beta-m(1-1/p)}(C_{\mathrm E}^dR^m)^{1-1/p}\\
 &\le C_{\mathrm{step}}^dN^{1/p}H R^\beta.
\end{align*}
The maximum-density proof of the same estimate observes pointwise that
\(\sum_wg_wF_w\le\max_wF_w\) and
\((\max_wF_w)^p\le\sum_wF_w^p\), then applies H\"older on
\(\mathcal E_{z,R}\).  It gives
\(\lambda_m(\mathcal E_{z,R})^{1-1/p}
(\sum_w\int F_w^p\,d\lambda_m)^{1/p}\), bounded by the same expression.
The radius exponents are
\(m(1-1/p)+\beta-m(1-1/p)=\beta\).
The densities in this calculation represent exact observations, so
the measurable routing functions are evaluated at the exact labels.
\end{proof}

The exact density can be replaced in this step by
Proposition~\ref{h:dv:measurable-test} and the test form of
Lemma~\ref{h:alt:selection}. For spherical restrictions,
Propositions~\ref{h:dv:joint-l1} and \ref{h:dv:pointwise-density}
give the same norm. Likewise, either enclosing-ball calculation can
replace the exact ellipsoid volume. These choices change only the
absolute constant, not the coefficient or radius powers.

\Needspace{4\baselineskip}
\begin{proposition}[A common finite success bound]\label{h:ev:coefficient-success}
Assume \(d\ge16\), set \(n=d-1\), \(p=\lfloor d/8\rfloor\), and let
\(n/2\le\beta\le d/2\). Let \(M,T\ge0\) be integers,
\(0<\epsilon\le1/10\), and consider a learner in
Definition~\ref{h:model}. Under the uniform-sphere signal prior, write
\(P_{\mathrm{succ}}\) for its success probability with either angular
or chordal error at most \(\epsilon\), and put \(b=\lceil T/p\rceil\).
Then
\begin{equation}\label{h:ev:common-success}
 P_{\mathrm{succ}}\le
 \min\left\{1,\,
 4^n\epsilon^{n-\beta}
 \left[C_{\mathrm{step}}^d
       \bigl((p+2)2^M\bigr)^{1/p}\right]^b\right\}.
\end{equation}
The same assertion holds with the initial factor \(4^n\) replaced
by the polar-cap constant \(C_0^d\).
\end{proposition}

\begin{proof}
Use the seedwise conditional experiment and Borel replacement from
Section~\ref{h:stream:reduction}. Conditional on almost every shared
seed, the initial state remains independent of the signal and rows.
Lemma~\ref{h:stream:categories} supplies, from every specified boundary
state, including an unreachable one, Borel routes to fresh
continuations. In a block of \(m\) rows their number is at most
\((m+2)2^M\): continuing states and terminal states tagged with one
of the \(m+1\) local stopping indices. The lemma applies to either
Borel terminal score.

The terminal bound \eqref{h:dv:terminal-coefficient} starts the
all-radius coefficient at \(H_0=4^n\epsilon^{n-\beta}\).
With \(N=(p+2)2^M\), Proposition~\ref{h:ev:coefficient-step} gives
one recurrence for all actual block lengths \(1\le m\le p\):
\begin{equation}\label{h:ev:coefficient-recursion}
 H_j=H_0\bigl(C_{\mathrm{step}}^dN^{1/p}\bigr)^j.
\end{equation}
Each multiplier is at least one, so terminal options, whose coefficient
is \(H_0\), satisfy every later bound. The ball \(B(0,1)\) contains
the whole sphere; hence the initial coefficient bounds total success.
Average first over the independent conditional initialization and then
over the original seed experiment. As in
Section~\ref{h:stream:reduction}, this averages original conditional
probabilities and requires no jointly measurable choice of the
seedwise replacements. The trivial probability bound gives the minimum
with one. For \(T=0\) this is the terminal estimate. Starting instead
from \eqref{h:ev:polar-terminal} proves the polar version.

One may instead pad the last block of actual length \(m<p\) with
independent Gaussian rows and extend each route by ignoring them.
Its matrix--label marginal on the first \(m\) rows is unchanged.
Thus full padded blocks give exactly the same estimate and count.
\end{proof}

The two endpoint exponents recover the following bounds.

\begin{corollary}[Finite success bounds for the two coefficients]
\label{h:ev:success}
Assume \(d\ge16\) and \(p=\lfloor d/8\rfloor\).
Let \(M,T\ge0\) be integers, let \(0<\epsilon\le1/10\), and consider
a learner in Definition~\ref{h:model} with these parameters.
Draw \(S\sim\sigma\) independently of its Gaussian rows
\(x_1,\ldots,x_T\) and its data-independent randomness, and set the
exact labels \(y_t=\langle x_t,S\rangle\).  Write \(\Pr_\sigma\) for
the resulting joint law and
put \(b=\lceil T/p\rceil\).  With an absolute
\(C_{\mathrm{step}}\ge1\), the uniform-prior angular success
probability satisfies
\begin{equation}\label{h:ev:half-sphere-success}
 \Pr_\sigma\{\operatorname{angle}(\widehat S,S)\le\epsilon\}
 \le
 \min\!\left\{1,\,
 4^d\epsilon^{(d-1)/2}
 \left[C_{\mathrm{step}}^d
       \bigl((p+2)2^M\bigr)^{1/p}\right]^b
 \right\}.
\end{equation}
The uniform-prior chordal success probability satisfies
\begin{equation}\label{h:ev:half-ambient-success}
 \Pr_\sigma\{\|\widehat S-S\|\le\epsilon\}
 \le
 \min\!\left\{1,\,
 C_0^d\epsilon^{d/2-1}
 \left[C_{\mathrm{step}}^d
       \bigl((2p+4)2^M\bigr)^{1/p}\right]^b
 \right\}.
\end{equation}
The second event contains the angular success event.
\end{corollary}

\begin{proof}
In Proposition~\ref{h:ev:coefficient-success}, set \(\beta=n/2\)
and use \(4^n\le4^d\) for the first formula. For the second, set
\(\beta=d/2\), use the polar terminal coefficient, and enlarge the
category bound to \(N=(2p+4)2^M\). The same recurrence then reads
\begin{equation}\label{h:ev:half-ambient-recursion}
 H_j=C_0^d\epsilon^{d/2-1}
       \bigl(C_{\mathrm{step}}^dN^{1/p}\bigr)^j.
\end{equation}
The final ball \(B(0,1)\) gives the second formula. Angular success
implies chordal success.
\end{proof}
